\clearpage\corrigeDevoir{4}
\begin{blocLivre}{dscor-4-1}\begin{corrigeBac}{6}{Chronophotographie}
\demarche{Utiliser les déplacements successifs et ne pas confondre les deux échelles.}
\schema[corrige]{chronophoto}
\begin{questions}
\item Trajectoire droite ; espacements croissants : valeur de vitesse croissante.\pts{2 pt}
\item $v_1=0,60/0,20=3,0$ m/s ; $v_3=1,40/0,20=7,0$ m/s.\pts{2 pt}
\item Longueurs 1,5 et 3,5 cm vers la droite. Le vecteur vitesse varie : les forces ne se compensent pas dans le référentiel assimilé galiléen.\pts{2 pt}
\end{questions}
\end{corrigeBac}\end{blocLivre}
\begin{blocLivre}{dscor-4-2}\begin{corrigeBac}{7}{Boule suspendue}
\demarche{Définir le système boule, puis utiliser le repos et le principe d’inertie.}
\schema[corrige]{suspension}
\begin{questions}
\item Terre sur boule : poids vers le bas ; fil sur boule : traction vers le haut.\pts{2 pt}
\item $P=mg=2,00$ N ; le repos donne $T=2,00$ N.\pts{2 pt}
\item Non : elles agissent toutes deux sur la boule.\pts{1 pt}
\item L’action gravitationnelle de la boule sur la Terre, de même valeur et de sens opposé.\pts{2 pt}
\end{questions}
\end{corrigeBac}\end{blocLivre}
\begin{blocLivre}{dscor-4-3}\begin{corrigeBac}{7}{Lancer vertical}
\demarche{Séparer vitesse à un instant et évolution de la vitesse sur une durée.}
\schema[corrige]{chute}
\begin{questions}
\item Chute libre : seule la gravitation, modélisée ici par le poids, agit.\pts{2 pt}
\item Poids vers le bas partout ; vitesse vers le haut, puis nulle, puis vers le bas.\pts{2 pt}
\item La vitesse change avant et après cet instant ; le poids reste non compensé.\pts{1 pt}
\item Non : une vitesse rectiligne constante suppose une compensation par des actions du liquide.\pts{2 pt}
\end{questions}
\end{corrigeBac}\end{blocLivre}

